% Options for packages loaded elsewhere
\PassOptionsToPackage{unicode}{hyperref}
\PassOptionsToPackage{hyphens}{url}
\documentclass[
  12pt,
]{ctexart}
\usepackage{xcolor}
\usepackage{amsmath,amssymb}
\setcounter{secnumdepth}{-\maxdimen} % remove section numbering
\usepackage{iftex}
\ifPDFTeX
  \usepackage[T1]{fontenc}
  \usepackage[utf8]{inputenc}
  \usepackage{textcomp} % provide euro and other symbols
\else % if luatex or xetex
  \usepackage{unicode-math} % this also loads fontspec
  \defaultfontfeatures{Scale=MatchLowercase}
  \defaultfontfeatures[\rmfamily]{Ligatures=TeX,Scale=1}
\fi
\usepackage{lmodern}
\ifPDFTeX\else
  % xetex/luatex font selection
\fi
% Use upquote if available, for straight quotes in verbatim environments
\IfFileExists{upquote.sty}{\usepackage{upquote}}{}
\IfFileExists{microtype.sty}{% use microtype if available
  \usepackage[]{microtype}
  \UseMicrotypeSet[protrusion]{basicmath} % disable protrusion for tt fonts
}{}
\makeatletter
\@ifundefined{KOMAClassName}{% if non-KOMA class
  \IfFileExists{parskip.sty}{%
    \usepackage{parskip}
  }{% else
    \setlength{\parindent}{0pt}
    \setlength{\parskip}{6pt plus 2pt minus 1pt}}
}{% if KOMA class
  \KOMAoptions{parskip=half}}
\makeatother
\usepackage{longtable,booktabs,array}
\usepackage{calc} % for calculating minipage widths
% Correct order of tables after \paragraph or \subparagraph
\usepackage{etoolbox}
\makeatletter
\patchcmd\longtable{\par}{\if@noskipsec\mbox{}\fi\par}{}{}
\makeatother
% Allow footnotes in longtable head/foot
\IfFileExists{footnotehyper.sty}{\usepackage{footnotehyper}}{\usepackage{footnote}}
\makesavenoteenv{longtable}
\setlength{\emergencystretch}{3em} % prevent overfull lines
\providecommand{\tightlist}{%
  \setlength{\itemsep}{0pt}\setlength{\parskip}{0pt}}
\usepackage{bookmark}
\IfFileExists{xurl.sty}{\usepackage{xurl}}{} % add URL line breaks if available
\urlstyle{same}
\hypersetup{
  hidelinks,
  pdfcreator={LaTeX via pandoc}}

\author{SCX}
\date{2026}

\begin{document}

\section{SCX TODO v4 ---
通用模块化重构}\label{scx-todo-v4-ux901aux7528ux6a21ux5757ux5316ux91cdux6784}

\begin{quote}
生成时间：2026-06-26
替代：\texttt{SCX\_TODO\_2026-06-26.md}（被本文件替代） 定位：综合
TODO，覆盖通用模块化重构 + 原等待项 详细规划：见
\texttt{SCX\_规划\_v4\_通用模块化.md}
\end{quote}

\begin{center}\rule{0.5\linewidth}{0.5pt}\end{center}

\subsection{标记说明}\label{ux6807ux8bb0ux8bf4ux660e}

\begin{longtable}[]{@{}lll@{}}
\toprule\noalign{}
标记 & 维度 & 说明 \\
\midrule\noalign{}
\endhead
\bottomrule\noalign{}
\endlastfoot
{[}A{]} & Academic / 学术 & 论文、投稿、学术影响力 \\
{[}C{]} & Code / 代码 & 软件工程、实现、开源 \\
{[}B{]} & Business / 商业 & 商业模式、合作、收入 \\
{[}IP{]} & Intellectual Property / 知识产权 & 专利、权属、保护 \\
{[}LLM{]} & LLM / 大模型 & 大模型数据评价方向 \\
\end{longtable}

\textbf{优先级}: P0=立即做, P1=1-2月内, P2=3-6月内, P3=6-12月

\textbf{时间标记}: ⚡=现在就能做, ⏳=需要等待前置条件

\begin{center}\rule{0.5\linewidth}{0.5pt}\end{center}

\subsection{Phase A: 核心重构 (2-4 周) 🔴
P0}\label{phase-a-ux6838ux5fc3ux91cdux6784-2-4-ux5468-p0}

\begin{quote}
目标：定义抽象接口，现有代码解耦，336 tests 通过
\end{quote}

\subsubsection{A.1
抽象基类定义}\label{a.1-ux62bdux8c61ux57faux7c7bux5b9aux4e49}

\begin{longtable}[]{@{}
  >{\raggedright\arraybackslash}p{(\linewidth - 12\tabcolsep) * \real{0.0682}}
  >{\raggedright\arraybackslash}p{(\linewidth - 12\tabcolsep) * \real{0.1364}}
  >{\raggedright\arraybackslash}p{(\linewidth - 12\tabcolsep) * \real{0.1818}}
  >{\raggedright\arraybackslash}p{(\linewidth - 12\tabcolsep) * \real{0.1364}}
  >{\raggedright\arraybackslash}p{(\linewidth - 12\tabcolsep) * \real{0.1364}}
  >{\raggedright\arraybackslash}p{(\linewidth - 12\tabcolsep) * \real{0.1364}}
  >{\raggedright\arraybackslash}p{(\linewidth - 12\tabcolsep) * \real{0.2045}}@{}}
\toprule\noalign{}
\begin{minipage}[b]{\linewidth}\raggedright
\#
\end{minipage} & \begin{minipage}[b]{\linewidth}\raggedright
TODO
\end{minipage} & \begin{minipage}[b]{\linewidth}\raggedright
优先级
\end{minipage} & \begin{minipage}[b]{\linewidth}\raggedright
依赖
\end{minipage} & \begin{minipage}[b]{\linewidth}\raggedright
预估
\end{minipage} & \begin{minipage}[b]{\linewidth}\raggedright
维度
\end{minipage} & \begin{minipage}[b]{\linewidth}\raggedright
可执行性
\end{minipage} \\
\midrule\noalign{}
\endhead
\bottomrule\noalign{}
\endlastfoot
A.1.1 & 定义 \texttt{SCXStateEncoder}
抽象基类（\texttt{featurize/distance/get\_state/state\_statistics}） &
P0 & 无 & 1d & {[}C{]} & ⚡ \\
A.1.2 & 定义 \texttt{SCXExpert}
抽象基类（\texttt{predict/confidence/cost/metadata}） & P0 & 无 & 0.5d &
{[}C{]} & ⚡ \\
A.1.3 & 定义 \texttt{SCXDataPoint}
抽象基类（\texttt{features/label/meta/state}） & P0 & 无 & 0.5d &
{[}C{]} & ⚡ \\
A.1.4 & 类型标注 + 文档字符串 + 抽象方法检查 & P0 & A.1.1-3 & 1d &
{[}C{]} & ⚡ \\
\end{longtable}

\subsubsection{A.2 Encoder 提取}\label{a.2-encoder-ux63d0ux53d6}

\begin{longtable}[]{@{}
  >{\raggedright\arraybackslash}p{(\linewidth - 12\tabcolsep) * \real{0.0682}}
  >{\raggedright\arraybackslash}p{(\linewidth - 12\tabcolsep) * \real{0.1364}}
  >{\raggedright\arraybackslash}p{(\linewidth - 12\tabcolsep) * \real{0.1818}}
  >{\raggedright\arraybackslash}p{(\linewidth - 12\tabcolsep) * \real{0.1364}}
  >{\raggedright\arraybackslash}p{(\linewidth - 12\tabcolsep) * \real{0.1364}}
  >{\raggedright\arraybackslash}p{(\linewidth - 12\tabcolsep) * \real{0.1364}}
  >{\raggedright\arraybackslash}p{(\linewidth - 12\tabcolsep) * \real{0.2045}}@{}}
\toprule\noalign{}
\begin{minipage}[b]{\linewidth}\raggedright
\#
\end{minipage} & \begin{minipage}[b]{\linewidth}\raggedright
TODO
\end{minipage} & \begin{minipage}[b]{\linewidth}\raggedright
优先级
\end{minipage} & \begin{minipage}[b]{\linewidth}\raggedright
依赖
\end{minipage} & \begin{minipage}[b]{\linewidth}\raggedright
预估
\end{minipage} & \begin{minipage}[b]{\linewidth}\raggedright
维度
\end{minipage} & \begin{minipage}[b]{\linewidth}\raggedright
可执行性
\end{minipage} \\
\midrule\noalign{}
\endhead
\bottomrule\noalign{}
\endlastfoot
A.2.1 & 从 \texttt{experiments/mlip\_case/} 提取
\texttt{MLIPEncoder}（ACE descriptor） & P0 & A.1.1 & 1d & {[}C{]} &
⚡ \\
A.2.2 & 从 \texttt{scx-health/} 提取
\texttt{VisionEncoder}（DINO/CLIP/ViT） & P0 & A.1.1 & 1d & {[}C{]} &
⚡ \\
A.2.3 & 从 \texttt{experiments/cifar/} 提取 \texttt{CIFAREncoder} & P0 &
A.1.1 & 0.5d & {[}C{]} & ⚡ \\
A.2.4 & 每个 encoder 独立模块，可单独 import & P0 & A.2.1-3 & 0.5d &
{[}C{]} & ⚡ \\
A.2.5 & 创建 \texttt{encoders/} 目录结构 & P0 & A.2.1-4 & 0.5d & {[}C{]}
& ⚡ \\
\end{longtable}

\subsubsection{A.3
核心模块解耦}\label{a.3-ux6838ux5fc3ux6a21ux5757ux89e3ux8026}

\begin{longtable}[]{@{}
  >{\raggedright\arraybackslash}p{(\linewidth - 12\tabcolsep) * \real{0.0682}}
  >{\raggedright\arraybackslash}p{(\linewidth - 12\tabcolsep) * \real{0.1364}}
  >{\raggedright\arraybackslash}p{(\linewidth - 12\tabcolsep) * \real{0.1818}}
  >{\raggedright\arraybackslash}p{(\linewidth - 12\tabcolsep) * \real{0.1364}}
  >{\raggedright\arraybackslash}p{(\linewidth - 12\tabcolsep) * \real{0.1364}}
  >{\raggedright\arraybackslash}p{(\linewidth - 12\tabcolsep) * \real{0.1364}}
  >{\raggedright\arraybackslash}p{(\linewidth - 12\tabcolsep) * \real{0.2045}}@{}}
\toprule\noalign{}
\begin{minipage}[b]{\linewidth}\raggedright
\#
\end{minipage} & \begin{minipage}[b]{\linewidth}\raggedright
TODO
\end{minipage} & \begin{minipage}[b]{\linewidth}\raggedright
优先级
\end{minipage} & \begin{minipage}[b]{\linewidth}\raggedright
依赖
\end{minipage} & \begin{minipage}[b]{\linewidth}\raggedright
预估
\end{minipage} & \begin{minipage}[b]{\linewidth}\raggedright
维度
\end{minipage} & \begin{minipage}[b]{\linewidth}\raggedright
可执行性
\end{minipage} \\
\midrule\noalign{}
\endhead
\bottomrule\noalign{}
\endlastfoot
A.3.1 & \texttt{scx/state/} 重构为使用 \texttt{SCXStateEncoder} 接口 &
P0 & A.1.1 & 2d & {[}C{]} & ⚡ \\
A.3.2 & \texttt{scx/expert/} 重构为使用 \texttt{SCXExpert} 接口 & P0 &
A.1.2 & 1d & {[}C{]} & ⚡ \\
A.3.3 & \texttt{scx/valuation/} 重构为使用统一 DataPoint 和 State 类型 &
P0 & A.1.3 & 1d & {[}C{]} & ⚡ \\
A.3.4 & \texttt{scx/action/} 重构为使用抽象接口 & P0 & A.1.1-3 & 1d &
{[}C{]} & ⚡ \\
A.3.5 & \texttt{scx/core/online.py} 使用抽象接口 & P0 & A.1.1-3 & 1d &
{[}C{]} & ⚡ \\
A.3.6 & 删除重复代码，统一数据流 & P0 & A.3.1-5 & 1d & {[}C{]} & ⚡ \\
\end{longtable}

\subsubsection{A.4
声明式配置原型}\label{a.4-ux58f0ux660eux5f0fux914dux7f6eux539fux578b}

\begin{longtable}[]{@{}
  >{\raggedright\arraybackslash}p{(\linewidth - 12\tabcolsep) * \real{0.0682}}
  >{\raggedright\arraybackslash}p{(\linewidth - 12\tabcolsep) * \real{0.1364}}
  >{\raggedright\arraybackslash}p{(\linewidth - 12\tabcolsep) * \real{0.1818}}
  >{\raggedright\arraybackslash}p{(\linewidth - 12\tabcolsep) * \real{0.1364}}
  >{\raggedright\arraybackslash}p{(\linewidth - 12\tabcolsep) * \real{0.1364}}
  >{\raggedright\arraybackslash}p{(\linewidth - 12\tabcolsep) * \real{0.1364}}
  >{\raggedright\arraybackslash}p{(\linewidth - 12\tabcolsep) * \real{0.2045}}@{}}
\toprule\noalign{}
\begin{minipage}[b]{\linewidth}\raggedright
\#
\end{minipage} & \begin{minipage}[b]{\linewidth}\raggedright
TODO
\end{minipage} & \begin{minipage}[b]{\linewidth}\raggedright
优先级
\end{minipage} & \begin{minipage}[b]{\linewidth}\raggedright
依赖
\end{minipage} & \begin{minipage}[b]{\linewidth}\raggedright
预估
\end{minipage} & \begin{minipage}[b]{\linewidth}\raggedright
维度
\end{minipage} & \begin{minipage}[b]{\linewidth}\raggedright
可执行性
\end{minipage} \\
\midrule\noalign{}
\endhead
\bottomrule\noalign{}
\endlastfoot
A.4.1 & 定义 YAML schema（encoder/expert/阈值/数据路径） & P0 & A.1.1-3
& 1d & {[}C{]} & ⚡ \\
A.4.2 & 实现 \texttt{ConfigLoader.from\_yaml(path)} → \texttt{SCXConfig}
& P0 & A.4.1 & 1d & {[}C{]} & ⚡ \\
A.4.3 & 创建 \texttt{domains/mlip.yaml} 示例配置 & P0 & A.4.1 & 0.5d &
{[}C{]} & ⚡ \\
A.4.4 & 创建 \texttt{domains/health.yaml} 示例配置 & P0 & A.4.1 & 0.5d &
{[}C{]} & ⚡ \\
A.4.5 & 创建 \texttt{domains/cifar.yaml} 示例配置 & P0 & A.4.1 & 0.5d &
{[}C{]} & ⚡ \\
\end{longtable}

\subsubsection{A.5 验证}\label{a.5-ux9a8cux8bc1}

\begin{longtable}[]{@{}
  >{\raggedright\arraybackslash}p{(\linewidth - 12\tabcolsep) * \real{0.0682}}
  >{\raggedright\arraybackslash}p{(\linewidth - 12\tabcolsep) * \real{0.1364}}
  >{\raggedright\arraybackslash}p{(\linewidth - 12\tabcolsep) * \real{0.1818}}
  >{\raggedright\arraybackslash}p{(\linewidth - 12\tabcolsep) * \real{0.1364}}
  >{\raggedright\arraybackslash}p{(\linewidth - 12\tabcolsep) * \real{0.1364}}
  >{\raggedright\arraybackslash}p{(\linewidth - 12\tabcolsep) * \real{0.1364}}
  >{\raggedright\arraybackslash}p{(\linewidth - 12\tabcolsep) * \real{0.2045}}@{}}
\toprule\noalign{}
\begin{minipage}[b]{\linewidth}\raggedright
\#
\end{minipage} & \begin{minipage}[b]{\linewidth}\raggedright
TODO
\end{minipage} & \begin{minipage}[b]{\linewidth}\raggedright
优先级
\end{minipage} & \begin{minipage}[b]{\linewidth}\raggedright
依赖
\end{minipage} & \begin{minipage}[b]{\linewidth}\raggedright
预估
\end{minipage} & \begin{minipage}[b]{\linewidth}\raggedright
维度
\end{minipage} & \begin{minipage}[b]{\linewidth}\raggedright
可执行性
\end{minipage} \\
\midrule\noalign{}
\endhead
\bottomrule\noalign{}
\endlastfoot
A.5.1 & 全部 336 tests 通过，无回归 & \textbf{P0} & A.2-4 & 1d & {[}C{]}
& ⚡ \\
A.5.2 & 每个领域 demo notebook 可运行 & P0 & A.2-4 & 1d & {[}C{]}{[}A{]}
& ⚡ \\
A.5.3 & 代码覆盖率不低于当前水平 & P0 & A.5.1 & 0.5d & {[}C{]} & ⚡ \\
\end{longtable}

\begin{center}\rule{0.5\linewidth}{0.5pt}\end{center}

\subsection{Phase B: 声明式配置 (2-3 周) 🟡
P1}\label{phase-b-ux58f0ux660eux5f0fux914dux7f6e-2-3-ux5468-p1}

\begin{quote}
目标：YAML → SCXFramework 自动构建，CLI 接口
\end{quote}

\begin{longtable}[]{@{}
  >{\raggedright\arraybackslash}p{(\linewidth - 12\tabcolsep) * \real{0.0682}}
  >{\raggedright\arraybackslash}p{(\linewidth - 12\tabcolsep) * \real{0.1364}}
  >{\raggedright\arraybackslash}p{(\linewidth - 12\tabcolsep) * \real{0.1818}}
  >{\raggedright\arraybackslash}p{(\linewidth - 12\tabcolsep) * \real{0.1364}}
  >{\raggedright\arraybackslash}p{(\linewidth - 12\tabcolsep) * \real{0.1364}}
  >{\raggedright\arraybackslash}p{(\linewidth - 12\tabcolsep) * \real{0.1364}}
  >{\raggedright\arraybackslash}p{(\linewidth - 12\tabcolsep) * \real{0.2045}}@{}}
\toprule\noalign{}
\begin{minipage}[b]{\linewidth}\raggedright
\#
\end{minipage} & \begin{minipage}[b]{\linewidth}\raggedright
TODO
\end{minipage} & \begin{minipage}[b]{\linewidth}\raggedright
优先级
\end{minipage} & \begin{minipage}[b]{\linewidth}\raggedright
依赖
\end{minipage} & \begin{minipage}[b]{\linewidth}\raggedright
预估
\end{minipage} & \begin{minipage}[b]{\linewidth}\raggedright
维度
\end{minipage} & \begin{minipage}[b]{\linewidth}\raggedright
可执行性
\end{minipage} \\
\midrule\noalign{}
\endhead
\bottomrule\noalign{}
\endlastfoot
B.1 & YAML 配置解析器完善（嵌套定义、路径映射、验证） & P1 & A.4 & 3d &
{[}C{]} & ⚡ \\
B.2 & \texttt{SCXFramework.from\_config(config)} 工厂方法 & P1 & A.3,
B.1 & 3d & {[}C{]} & ⚡ \\
B.3 & CLI
设计：\texttt{scx\ run/config-check/list-encoders/list-experts} & P1 &
B.2 & 2d & {[}C{]} & ⚡ \\
B.4 & 审计报告自动生成（Markdown/PDF/HTML + 可视化） & P1 & B.2 & 3d &
{[}C{]}{[}B{]} & ⚡ \\
B.5 & Encoder 开发指南 + YAML 配置参考 + 端到端 demo & P1 & B.2-4 & 2d &
{[}C{]}{[}A{]} & ⚡ \\
\end{longtable}

\begin{center}\rule{0.5\linewidth}{0.5pt}\end{center}

\subsection{Phase C: 新领域扩展 (4-8 周) 🟡
P1}\label{phase-c-ux65b0ux9886ux57dfux6269ux5c55-4-8-ux5468-p1}

\begin{quote}
目标：用新架构快速适配新领域，验证通用性
\end{quote}

\subsubsection{C.1 SCX-Robot}\label{c.1-scx-robot}

\begin{longtable}[]{@{}
  >{\raggedright\arraybackslash}p{(\linewidth - 12\tabcolsep) * \real{0.0682}}
  >{\raggedright\arraybackslash}p{(\linewidth - 12\tabcolsep) * \real{0.1364}}
  >{\raggedright\arraybackslash}p{(\linewidth - 12\tabcolsep) * \real{0.1818}}
  >{\raggedright\arraybackslash}p{(\linewidth - 12\tabcolsep) * \real{0.1364}}
  >{\raggedright\arraybackslash}p{(\linewidth - 12\tabcolsep) * \real{0.1364}}
  >{\raggedright\arraybackslash}p{(\linewidth - 12\tabcolsep) * \real{0.1364}}
  >{\raggedright\arraybackslash}p{(\linewidth - 12\tabcolsep) * \real{0.2045}}@{}}
\toprule\noalign{}
\begin{minipage}[b]{\linewidth}\raggedright
\#
\end{minipage} & \begin{minipage}[b]{\linewidth}\raggedright
TODO
\end{minipage} & \begin{minipage}[b]{\linewidth}\raggedright
优先级
\end{minipage} & \begin{minipage}[b]{\linewidth}\raggedright
依赖
\end{minipage} & \begin{minipage}[b]{\linewidth}\raggedright
预估
\end{minipage} & \begin{minipage}[b]{\linewidth}\raggedright
维度
\end{minipage} & \begin{minipage}[b]{\linewidth}\raggedright
可执行性
\end{minipage} \\
\midrule\noalign{}
\endhead
\bottomrule\noalign{}
\endlastfoot
C.1.1 & 调研 LeRobot 公开数据 / Open X-Embodiment 子集 & P1 & 无 & 2d &
{[}A{]}{[}C{]} & ⚡ \\
C.1.2 & 实现 \texttt{TrajectoryEncoder}（视觉+关节+力+成功信号融合） &
P1 & A.1.1 & 5d & {[}C{]} & ⚡ \\
C.1.3 & 创建 \texttt{domains/robot.yaml} & P1 & B.1 & 0.5d & {[}C{]} &
⚡ \\
C.1.4 & 成功/失败状态分类实验 & P1 & C.1.2 & 3d & {[}A{]}{[}C{]} & ⏳等
GPU \\
C.1.5 & 冗余轨迹压缩实验 & P1 & C.1.2 & 2d & {[}A{]}{[}C{]} & ⏳等
GPU \\
C.1.6 & 失败状态挖掘实验 & P1 & C.1.2 & 2d & {[}A{]}{[}C{]} & ⏳等
GPU \\
C.1.7 & 对比：full-data vs SCX-selected vs random 在 imitation policy
上表现 & P1 & C.1.4-6 & 3d & {[}A{]} & ⏳等 GPU \\
\end{longtable}

\subsubsection{C.2 SCX-FEM}\label{c.2-scx-fem}

\begin{longtable}[]{@{}
  >{\raggedright\arraybackslash}p{(\linewidth - 12\tabcolsep) * \real{0.0682}}
  >{\raggedright\arraybackslash}p{(\linewidth - 12\tabcolsep) * \real{0.1364}}
  >{\raggedright\arraybackslash}p{(\linewidth - 12\tabcolsep) * \real{0.1818}}
  >{\raggedright\arraybackslash}p{(\linewidth - 12\tabcolsep) * \real{0.1364}}
  >{\raggedright\arraybackslash}p{(\linewidth - 12\tabcolsep) * \real{0.1364}}
  >{\raggedright\arraybackslash}p{(\linewidth - 12\tabcolsep) * \real{0.1364}}
  >{\raggedright\arraybackslash}p{(\linewidth - 12\tabcolsep) * \real{0.2045}}@{}}
\toprule\noalign{}
\begin{minipage}[b]{\linewidth}\raggedright
\#
\end{minipage} & \begin{minipage}[b]{\linewidth}\raggedright
TODO
\end{minipage} & \begin{minipage}[b]{\linewidth}\raggedright
优先级
\end{minipage} & \begin{minipage}[b]{\linewidth}\raggedright
依赖
\end{minipage} & \begin{minipage}[b]{\linewidth}\raggedright
预估
\end{minipage} & \begin{minipage}[b]{\linewidth}\raggedright
维度
\end{minipage} & \begin{minipage}[b]{\linewidth}\raggedright
可执行性
\end{minipage} \\
\midrule\noalign{}
\endhead
\bottomrule\noalign{}
\endlastfoot
C.2.1 & 调研 PDEBench / 自生成弹性/热传导数据 & P1 & 无 & 2d &
{[}A{]}{[}C{]} & ⚡ \\
C.2.2 & 实现 \texttt{SimulationEncoder}（网格/几何/响应 embedding） & P1
& A.1.1 & 5d & {[}C{]} & ⚡ \\
C.2.3 & 创建 \texttt{domains/fem.yaml} & P1 & B.1 & 0.5d & {[}C{]} &
⚡ \\
C.2.4 & 多保真调度实验：粗网格 vs 细网格 & P1 & C.2.2 & 3d &
{[}A{]}{[}C{]} & ⚡ \\
C.2.5 & Surrogate 可靠性判断实验 & P1 & C.2.2 & 2d & {[}A{]}{[}C{]} &
⚡ \\
C.2.6 & 对比：SCX 多保真调度 vs 全细网格 vs 全粗网格 & P1 & C.2.4-5 & 2d
& {[}A{]} & ⏳等 GPU \\
\end{longtable}

\subsubsection{C.3 SCX-LLM (minimal)}\label{c.3-scx-llm-minimal}

\begin{longtable}[]{@{}
  >{\raggedright\arraybackslash}p{(\linewidth - 12\tabcolsep) * \real{0.0682}}
  >{\raggedright\arraybackslash}p{(\linewidth - 12\tabcolsep) * \real{0.1364}}
  >{\raggedright\arraybackslash}p{(\linewidth - 12\tabcolsep) * \real{0.1818}}
  >{\raggedright\arraybackslash}p{(\linewidth - 12\tabcolsep) * \real{0.1364}}
  >{\raggedright\arraybackslash}p{(\linewidth - 12\tabcolsep) * \real{0.1364}}
  >{\raggedright\arraybackslash}p{(\linewidth - 12\tabcolsep) * \real{0.1364}}
  >{\raggedright\arraybackslash}p{(\linewidth - 12\tabcolsep) * \real{0.2045}}@{}}
\toprule\noalign{}
\begin{minipage}[b]{\linewidth}\raggedright
\#
\end{minipage} & \begin{minipage}[b]{\linewidth}\raggedright
TODO
\end{minipage} & \begin{minipage}[b]{\linewidth}\raggedright
优先级
\end{minipage} & \begin{minipage}[b]{\linewidth}\raggedright
依赖
\end{minipage} & \begin{minipage}[b]{\linewidth}\raggedright
预估
\end{minipage} & \begin{minipage}[b]{\linewidth}\raggedright
维度
\end{minipage} & \begin{minipage}[b]{\linewidth}\raggedright
可执行性
\end{minipage} \\
\midrule\noalign{}
\endhead
\bottomrule\noalign{}
\endlastfoot
C.3.1 & 实现 \texttt{TextEncoder}（LLM embedding + 任务类型编码） & P1 &
A.1.1 & 3d & {[}C{]} & ⚡ \\
C.3.2 & 创建 \texttt{domains/llm.yaml} (minimal) & P1 & B.1 & 0.5d &
{[}C{]} & ⚡ \\
C.3.3 & 极小验证：区分有价值/冗余/噪声的 instruction 数据 & P2 & C.3.1 &
3d & {[}A{]}{[}LLM{]} & ⏳等 GPU \\
C.3.4 & 定位文档：SCX-LLM 聚焦可验证任务（数学/代码/逻辑） & P2 & 无 &
2d & {[}LLM{]}{[}A{]} & ⚡ \\
\end{longtable}

\begin{center}\rule{0.5\linewidth}{0.5pt}\end{center}

\subsection{Phase D: 插件系统 (长期) 🟢
P2}\label{phase-d-ux63d2ux4ef6ux7cfbux7edf-ux957fux671f-p2}

\begin{quote}
目标：模块插件化，社区贡献指南，Benchmark 标准化
\end{quote}

\begin{longtable}[]{@{}
  >{\raggedright\arraybackslash}p{(\linewidth - 12\tabcolsep) * \real{0.0682}}
  >{\raggedright\arraybackslash}p{(\linewidth - 12\tabcolsep) * \real{0.1364}}
  >{\raggedright\arraybackslash}p{(\linewidth - 12\tabcolsep) * \real{0.1818}}
  >{\raggedright\arraybackslash}p{(\linewidth - 12\tabcolsep) * \real{0.1364}}
  >{\raggedright\arraybackslash}p{(\linewidth - 12\tabcolsep) * \real{0.1364}}
  >{\raggedright\arraybackslash}p{(\linewidth - 12\tabcolsep) * \real{0.1364}}
  >{\raggedright\arraybackslash}p{(\linewidth - 12\tabcolsep) * \real{0.2045}}@{}}
\toprule\noalign{}
\begin{minipage}[b]{\linewidth}\raggedright
\#
\end{minipage} & \begin{minipage}[b]{\linewidth}\raggedright
TODO
\end{minipage} & \begin{minipage}[b]{\linewidth}\raggedright
优先级
\end{minipage} & \begin{minipage}[b]{\linewidth}\raggedright
依赖
\end{minipage} & \begin{minipage}[b]{\linewidth}\raggedright
预估
\end{minipage} & \begin{minipage}[b]{\linewidth}\raggedright
维度
\end{minipage} & \begin{minipage}[b]{\linewidth}\raggedright
可执行性
\end{minipage} \\
\midrule\noalign{}
\endhead
\bottomrule\noalign{}
\endlastfoot
D.1 & 插件注册机制（\texttt{register()} 装饰器 + 自动发现 + 沙箱） & P2
& A.3 & 1w & {[}C{]} & ⚡ \\
D.2 & \texttt{SCXCompress} 插件化 & P2 & D.1 & 2d & {[}C{]} & ⚡ \\
D.3 & \texttt{OnlineSCXFramework} 插件化 & P2 & D.1 & 2d & {[}C{]} &
⚡ \\
D.4 & \texttt{StateConditionedInfluence} 插件化 & P2 & D.1 & 1d &
{[}C{]} & ⚡ \\
D.5 & CONTRIBUTING.md + CLA + Issue/PR 模板 & P2 & D.1-4 & 1w &
{[}C{]}{[}B{]} & ⚡ \\
D.6 & SCX-Bench 标准化 API + 5-10 数据集 + Baseline 集成 & P2 & D.1 & 2w
& {[}C{]}{[}A{]} & ⏳等 GPU \\
\end{longtable}

\begin{center}\rule{0.5\linewidth}{0.5pt}\end{center}

\subsection{原等待项（保留，来自 Phase
1-3）}\label{ux539fux7b49ux5f85ux9879ux4fddux7559ux6765ux81ea-phase-1-3}

以下为之前 TODO 中的等待项，与本 TODO 并行推进。

\subsubsection{等待 GPU
(8月中旬，出差回来)}\label{ux7b49ux5f85-gpu-8ux6708ux4e2dux65ecux51faux5deeux56deux6765}

\begin{longtable}[]{@{}
  >{\raggedright\arraybackslash}p{(\linewidth - 10\tabcolsep) * \real{0.0732}}
  >{\raggedright\arraybackslash}p{(\linewidth - 10\tabcolsep) * \real{0.1463}}
  >{\raggedright\arraybackslash}p{(\linewidth - 10\tabcolsep) * \real{0.1951}}
  >{\raggedright\arraybackslash}p{(\linewidth - 10\tabcolsep) * \real{0.2195}}
  >{\raggedright\arraybackslash}p{(\linewidth - 10\tabcolsep) * \real{0.1463}}
  >{\raggedright\arraybackslash}p{(\linewidth - 10\tabcolsep) * \real{0.2195}}@{}}
\toprule\noalign{}
\begin{minipage}[b]{\linewidth}\raggedright
\#
\end{minipage} & \begin{minipage}[b]{\linewidth}\raggedright
TODO
\end{minipage} & \begin{minipage}[b]{\linewidth}\raggedright
优先级
\end{minipage} & \begin{minipage}[b]{\linewidth}\raggedright
阻塞因素
\end{minipage} & \begin{minipage}[b]{\linewidth}\raggedright
维度
\end{minipage} & \begin{minipage}[b]{\linewidth}\raggedright
预计解锁
\end{minipage} \\
\midrule\noalign{}
\endhead
\bottomrule\noalign{}
\endlastfoot
W.1 & CIFAR/MedMNIST GPU 重跑（ResNet-18 + 50 epochs） & P1 & 无 GPU &
{[}A{]}{[}C{]} & 8月中旬 \\
W.2 & SCX-Noise 论文 Figure (publication-quality) & P1 & W.1 & {[}A{]} &
8月中旬 \\
W.3 & 与 Data Shapley/LESS 对比实验（GPU加速后） & P1 & W.1 &
{[}A{]}{[}C{]} & 8月中旬 \\
W.4 & HAM10000/ISIC 实验 & P1 & W.1 & {[}A{]}{[}C{]} & 8月中旬 \\
W.5 & SCX-Health 论文草稿 & P2 & W.4 & {[}A{]} & 8月中旬 \\
W.6 & SCX-Robot 实验 (C.1.4-7) & P1 & W.1 & {[}A{]}{[}C{]} & 8月中旬 \\
W.7 & SCX-FEM GPU 实验 (C.2.6) & P1 & W.1 & {[}A{]} & 8月中旬 \\
W.8 & SCX-LLM 实验 (C.3.3) & P2 & W.1 & {[}LLM{]}{[}A{]} & 8月中旬 \\
W.9 & SCX-Bench GPU baseline 集成 (D.6) & P2 & W.1 & {[}C{]}{[}A{]} &
8月中旬 \\
\end{longtable}

\subsubsection{等待 DFT
(2-4周，超算运行中)}\label{ux7b49ux5f85-dft-2-4ux5468ux8d85ux7b97ux8fd0ux884cux4e2d}

\begin{longtable}[]{@{}
  >{\raggedright\arraybackslash}p{(\linewidth - 10\tabcolsep) * \real{0.0732}}
  >{\raggedright\arraybackslash}p{(\linewidth - 10\tabcolsep) * \real{0.1463}}
  >{\raggedright\arraybackslash}p{(\linewidth - 10\tabcolsep) * \real{0.1951}}
  >{\raggedright\arraybackslash}p{(\linewidth - 10\tabcolsep) * \real{0.2195}}
  >{\raggedright\arraybackslash}p{(\linewidth - 10\tabcolsep) * \real{0.1463}}
  >{\raggedright\arraybackslash}p{(\linewidth - 10\tabcolsep) * \real{0.2195}}@{}}
\toprule\noalign{}
\begin{minipage}[b]{\linewidth}\raggedright
\#
\end{minipage} & \begin{minipage}[b]{\linewidth}\raggedright
TODO
\end{minipage} & \begin{minipage}[b]{\linewidth}\raggedright
优先级
\end{minipage} & \begin{minipage}[b]{\linewidth}\raggedright
阻塞因素
\end{minipage} & \begin{minipage}[b]{\linewidth}\raggedright
维度
\end{minipage} & \begin{minipage}[b]{\linewidth}\raggedright
预计解锁
\end{minipage} \\
\midrule\noalign{}
\endhead
\bottomrule\noalign{}
\endlastfoot
W.10 & GaN\_v1 DFT 结果 → 训练 GaN expert & P1 & 超算运行中 &
{[}A{]}{[}C{]} & 2-4周 \\
W.11 & AlN\_v4 DFT 结果 → 重训 AlN expert & P1 & 超算运行中 &
{[}A{]}{[}C{]} & 2-4周 \\
W.12 & EGP Paper 1 完整实验 (M1-M6) & P1 & W.10-11 & {[}A{]} & 2-4周 \\
W.13 & SCX-MLIP 论文完成 & P1 & W.10-11 & {[}A{]} & 2-4周 \\
W.14 & SCX-Potential Compiler 实验 (2.11-2.18) & P2 & W.10-11 &
{[}C{]}{[}A{]} & 2-4周 \\
\end{longtable}

\subsubsection{论文与投稿}\label{ux8bbaux6587ux4e0eux6295ux7a3f}

\begin{longtable}[]{@{}
  >{\raggedright\arraybackslash}p{(\linewidth - 10\tabcolsep) * \real{0.0789}}
  >{\raggedright\arraybackslash}p{(\linewidth - 10\tabcolsep) * \real{0.1579}}
  >{\raggedright\arraybackslash}p{(\linewidth - 10\tabcolsep) * \real{0.2105}}
  >{\raggedright\arraybackslash}p{(\linewidth - 10\tabcolsep) * \real{0.1579}}
  >{\raggedright\arraybackslash}p{(\linewidth - 10\tabcolsep) * \real{0.1579}}
  >{\raggedright\arraybackslash}p{(\linewidth - 10\tabcolsep) * \real{0.2368}}@{}}
\toprule\noalign{}
\begin{minipage}[b]{\linewidth}\raggedright
\#
\end{minipage} & \begin{minipage}[b]{\linewidth}\raggedright
TODO
\end{minipage} & \begin{minipage}[b]{\linewidth}\raggedright
优先级
\end{minipage} & \begin{minipage}[b]{\linewidth}\raggedright
依赖
\end{minipage} & \begin{minipage}[b]{\linewidth}\raggedright
维度
\end{minipage} & \begin{minipage}[b]{\linewidth}\raggedright
可执行性
\end{minipage} \\
\midrule\noalign{}
\endhead
\bottomrule\noalign{}
\endlastfoot
W.15 & EGP Paper 1 Introduction + Methods 写作 & \textbf{P0} & 无 &
{[}A{]} & ⚡ \\
W.16 & SCX-Theory arXiv 草稿 & \textbf{P0} & 无 & {[}A{]} & ⚡ \\
W.17 & 投稿 npj Computational Materials / Nature Communications & P1 &
W.13 & {[}A{]} & ⏳ \\
W.18 & arXiv 占坑（SCX-Theory） & P1 & W.16 & {[}A{]} & ⚡ \\
W.19 & SCX-Sim 大论文 & P2 & W.17 & {[}A{]} & ⏳ \\
W.20 & SCX-Health 论文投稿 & P2 & W.5 & {[}A{]} & ⏳ \\
\end{longtable}

\subsubsection{商业}\label{ux5546ux4e1a}

\begin{longtable}[]{@{}llllll@{}}
\toprule\noalign{}
\# & TODO & 优先级 & 依赖 & 维度 & 可执行性 \\
\midrule\noalign{}
\endhead
\bottomrule\noalign{}
\endlastfoot
W.21 & 调研华为健康合作入口 & P1 & 无 & {[}B{]} & ⚡ \\
W.22 & 调研国产 EDA/TCAD 公司技术合作 & P2 & 无 & {[}B{]} & ⚡ \\
W.23 & 找 1-2 个 paid pilot 甲方 & P2 & 有 demo 后 & {[}B{]} & ⏳ \\
W.24 & 设计三层商业产品线 & P2 & 无 & {[}B{]}{[}C{]} & ⚡ \\
W.25 & 找校外 IP 律师评估 & P1 & 无 & {[}IP{]} & ⚡ \\
\end{longtable}

\subsubsection{IP 与权属}\label{ip-ux4e0eux6743ux5c5e}

\begin{longtable}[]{@{}llllll@{}}
\toprule\noalign{}
\# & TODO & 优先级 & 依赖 & 维度 & 可执行性 \\
\midrule\noalign{}
\endhead
\bottomrule\noalign{}
\endlastfoot
W.26 & 关键代码 Sha-256 存证 & P1 & 无 & {[}IP{]} & ⚡ \\
W.27 & 专利价值判断 & P1 & W.25 & {[}IP{]} & ⏳ \\
W.28 & CLA 模板 & P2 & 无 & {[}IP{]} & ⚡ \\
\end{longtable}

\begin{center}\rule{0.5\linewidth}{0.5pt}\end{center}

\subsection{时间线总览}\label{ux65f6ux95f4ux7ebfux603bux89c8}

\begin{verbatim}
现在 (6月)                   7月                        8月                        9月
─────                       ───                        ───                        ───

Phase A ───────────────────→ Phase B ─────────────────→ Phase C ────────────────→ Phase D
核心重构   2-4周             声明式配置  2-3周            新领域扩展  4-8周          插件系统
                                                                                 
  A.1 抽象基类                 B.1 YAML 完善                C.1 SCX-Robot           D.1 插件注册
  A.2 Encoder提取             B.2 Factory 方法             C.2 SCX-FEM              D.2-4 插件化
  A.3 核心解耦                B.3 CLI                     C.3 SCX-LLM (minimal)     D.5 社区指南
  A.4 配置原型                B.4 报告生成                                          D.6 Benchmark
  A.5 验证                    B.5 文档
                                                                                 
等待项 ──────────────────────────────────────────────────────────────────────────
                                                                                 
W.15 EGP Paper 1 写作 ──────→ W.16 SCX-Theory arXiv ────→ W.13 MLIP 论文完成 ────→ W.17 投稿
                              W.18 arXiv 占坑                                    
                              W.10-12 DFT 结果 ↔ W.14 Pot Compiler
                              W.1-9 GPU 实验 (8月中旬) ──→ W.4-5 Health 论文
                                                          W.6-8 Robot/FEM/LLM
\end{verbatim}

\begin{center}\rule{0.5\linewidth}{0.5pt}\end{center}

\subsection{依赖关系图}\label{ux4f9dux8d56ux5173ux7cfbux56fe}

\begin{verbatim}
Phase A (核心重构)             Phase B (声明式配置)         Phase C (新领域)
─────────────────             ────────────────────        ─────────────────

A.1 抽象基类 ────→ A.2 Encoder提取
                      │
                      ├──→ A.3 核心解耦 ────→ A.5 验证
                      │                        │
                      └──→ A.4 配置原型 ────→ B.1 YAML完善 → B.2 Factory → B.3 CLI
                                                   │            │
                                                   └──→ B.4 报告生成
                                                        B.5 文档
                                                        │
                              C.1 Robot ────────────────┘
                              C.2 FEM  ─────────────────┘
                              C.3 LLM  ─────────────────┘

等待项:

W.15 EGP Paper 1 ───→ W.16 SCX-Theory arXiv ───→ W.18 arXiv 占坑
W.10 DFT结果   ───→ W.12 EGP实验     ───→ W.13 MLIP论文
W.1 GPU       ───→ W.2-9 全部GPU实验
\end{verbatim}

\begin{center}\rule{0.5\linewidth}{0.5pt}\end{center}

\subsection{核心风险提示}\label{ux6838ux5fc3ux98ceux9669ux63d0ux793a}

\begin{longtable}[]{@{}
  >{\raggedright\arraybackslash}p{(\linewidth - 4\tabcolsep) * \real{0.2500}}
  >{\raggedright\arraybackslash}p{(\linewidth - 4\tabcolsep) * \real{0.3750}}
  >{\raggedright\arraybackslash}p{(\linewidth - 4\tabcolsep) * \real{0.3750}}@{}}
\toprule\noalign{}
\begin{minipage}[b]{\linewidth}\raggedright
风险
\end{minipage} & \begin{minipage}[b]{\linewidth}\raggedright
严重程度
\end{minipage} & \begin{minipage}[b]{\linewidth}\raggedright
缓解措施
\end{minipage} \\
\midrule\noalign{}
\endhead
\bottomrule\noalign{}
\endlastfoot
重构破坏现有 336 tests & 🔴 高 & 每步重构保持 test 覆盖，CI pipeline \\
过度抽象导致灵活性下降 & 🟡 中 & 保留 \texttt{CustomEncoder} 逃生口 \\
新领域 encoder 开发成本超预期 & 🟡 中 & 先用简单 embedding 做 MVP \\
学校主张 SCX IP 权属 & 🔴 高 & Phase 0 证据隔离 + 校外律师评估 \\
GPU 到位后实验时间紧 & 🟡 中 & 提前准备好代码，GPU 一到就开跑 \\
论文审稿人质疑''框架论文'' & 🟡 中 &
论文重心在具体实验，不单独发框架论文 \\
VASP 任务占满超算 & 🟢 低 & 分批提交、监控负载 \\
\end{longtable}

\begin{center}\rule{0.5\linewidth}{0.5pt}\end{center}

\subsection{总结：现在立刻做的 5
件事}\label{ux603bux7ed3ux73b0ux5728ux7acbux523bux505aux7684-5-ux4ef6ux4e8b}

\begin{enumerate}
\def\labelenumi{\arabic{enumi}.}
\tightlist
\item
  \textbf{定义抽象基类} (A.1.1-A.1.4) --- 通用架构的根基
\item
  \textbf{提取现有 encoder} (A.2.1-A.2.5) --- 解耦第一步
\item
  \textbf{重构核心模块} (A.3.1-A.3.6) --- 让代码只依赖抽象接口
\item
  \textbf{创建 YAML 配置原型} (A.4.1-A.4.5) --- 声明式配置的第一步
\item
  \textbf{保持 W.15/W.16 论文写作并行} --- 架构重构不影响论文进度
\end{enumerate}

\begin{quote}
\textbf{一句话战略：先解耦再扩展。Phase A
完成后，新增领域的适配成本将从''2-3 周''降到''2-3 天''。}
\end{quote}

\end{document}
